\documentclass[11pt]{article}

\usepackage[T1]{fontenc}
\usepackage[utf8]{inputenc}
\usepackage{lmodern}
\usepackage{amsmath,amssymb,amsfonts,amsthm}
\usepackage[margin=1in]{geometry}
\usepackage{microtype}
\usepackage{xcolor}
\usepackage[colorlinks=true,linkcolor=teal,urlcolor=teal]{hyperref}

\newcommand{\F}{\mathcal F}
\newcommand{\Hc}{\mathcal H}
\newcommand{\T}{\mathcal T_1}
\newcommand{\Bop}{\mathcal B}
\newcommand{\C}{\mathbb C}
\newcommand{\Tr}{\operatorname{Tr}}

\newtheorem*{lemmaD}{Lemma D.1 (rephrased)}

\title{A Short Guide to Lemma D.1\\
\large Why a covariant quantum operation has constant success probability}
\author{}
\date{}

\begin{document}
\maketitle

\begin{abstract}
Lemma D.1 says that a displacement-covariant, trace-nonincreasing quantum
operation can lose only a fixed fraction of the trace: its success
probability cannot depend on the input state.  The proof is a short
operator-theoretic version of Schur's lemma.  This note explains the duality
behind the proof, the irreducibility of the Weyl representation, and the
lemma's role in the Gaussian stochastic-order argument.
\end{abstract}

\section{The content in one sentence}

Let $\Gamma$ be a possibly probabilistic operation on one or finitely many
oscillator modes.  If displacing the input by $z$ merely displaces the output
by $\sqrt\gamma z$, then no displacement can change the probability that the
operation succeeds.  Because the Weyl displacement representation is
irreducible, this invariance is strong enough to make the success probability
the same for \emph{every} input state, not only for states on one displacement
orbit.

Equivalently, there is a single number $c\in[0,1]$ such that
\[
  \Tr\Gamma(\rho)=c
\]
for every normalized state $\rho$.  When $c>0$, the rescaled map
$c^{-1}\Gamma$ is a deterministic channel.

\section{Background}

\subsection{Trace class and bounded operators}

For $s<\infty$ oscillator modes, let
\[
  \Hc=\F^{\otimes s},
  \qquad
  \T(\Hc)=\{\text{trace-class operators on }\Hc\},
  \qquad
  \Bop(\Hc)=\{\text{bounded operators on }\Hc\}.
\]
The two spaces are in duality:
\[
  \T(\Hc)^*=\Bop(\Hc),
  \qquad
  \langle X,A\rangle=\Tr(XA).
\]
Thus every bounded linear functional $\ell$ on the trace class has a unique
representing operator $A\in\Bop(\Hc)$:
\begin{equation}\label{eq:dual-representation}
  \ell(X)=\Tr(AX).
\end{equation}
Moreover, $\ell$ is positive exactly when $A\geq0$.

If $\Gamma:\T(\Hc)\to\T(\Hc)$ is bounded, its Banach adjoint is the map
\[
  \Gamma^*:\Bop(\Hc)\longrightarrow\Bop(\Hc),
  \qquad
  \Tr[\Gamma(X)A]=\Tr[X\Gamma^*(A)].
\]
In particular, the functional recording the output trace is represented by
the \emph{effect}
\begin{equation}\label{eq:effect}
  A=\Gamma^*(I),
  \qquad
  \Tr\Gamma(X)=\Tr(AX).
\end{equation}
If $\Gamma$ is completely positive and trace nonincreasing, then
\[
  0\leq A\leq I.
\]
Indeed, positivity gives $A\geq0$, while
$\Tr[\Gamma(X)]\leq\Tr X$ for every $X\geq0$ gives $A\leq I$ by duality.
For a normalized input $\rho$, $\Tr(A\rho)$ is precisely the success
probability of the operation.

\subsection{Weyl displacements}

On one mode,
\[
  D(z)=\exp(za^*-\bar z a),\qquad z\in\C.
\]
On $s$ modes, $z=(z_1,\ldots,z_s)\in\C^s$ and
\[
  D(z)=\bigotimes_{j=1}^sD(z_j).
\]
The Weyl relation contains a scalar phase,
\[
  D(z)D(w)=e^{\mathrm i\operatorname{Im}(z\cdot\bar w)}D(z+w),
\]
but this phase disappears under conjugation.  Consequently,
$X\mapsto D(z)XD(z)^*$ is an ordinary action of the additive phase-space
group.

The Weyl representation on the $s$-mode Fock space is irreducible.  Its
commutant is therefore
\begin{equation}\label{eq:weyl-commutant}
  \{A\in\Bop(\Hc):AD(z)=D(z)A\text{ for every }z\in\C^s\}
  =\C I.
\end{equation}
This is the Stone--von Neumann irreducibility statement, followed by Schur's
lemma.  Informally, commuting with all position and momentum translations
leaves an operator no nontrivial degree of freedom.

\subsection{Displacement covariance}

The covariance condition in the paper is
\begin{equation}\label{eq:covariance}
  \Gamma\bigl(D(z)XD(z)^*\bigr)
  =D(\sqrt\gamma z)\Gamma(X)D(\sqrt\gamma z)^*.
\end{equation}
The factor $\sqrt\gamma$ is the amplitude gain.  Its value will not enter the
proof: only the fact that the output transformation is unitary, and hence
trace preserving, is needed.

\section{Statement of Lemma D.1}

\begin{lemmaD}
Let $\Gamma:\T(\Hc)\to\T(\Hc)$ be completely positive and trace
nonincreasing, and suppose that \eqref{eq:covariance} holds for every
$z\in\C^s$.  Then there is a constant $c\in[0,1]$ such that
\begin{equation}\label{eq:constant-trace}
  \Tr\Gamma(X)=c\,\Tr X
\end{equation}
for every trace-class operator $X$.
\end{lemmaD}

The statement includes subnormalized and nonpositive $X$.  It is enough to
prove the identity first for positive $X$; linearity then extends it to all
trace-class operators.

\section{Proof, line by line}

\subsection*{Step 1: encode the output trace by one operator}

Let $A=\Gamma^*(I)$.  By \eqref{eq:effect},
\[
  \Tr\Gamma(X)=\Tr(AX),
\]
and trace nonincrease gives $0\leq A\leq I$.  Thus all possible
input-dependence of the success probability is stored in the single effect
$A$.

\subsection*{Step 2: covariance makes the effect invariant}

Take the trace of both sides of \eqref{eq:covariance}.  Unitary invariance of
the trace gives
\[
  \Tr\Gamma(D(z)XD(z)^*)=\Tr\Gamma(X).
\]
Using the representation by $A$ and cyclicity of the trace,
\begin{align*}
  \Tr\bigl[A D(z)XD(z)^*\bigr]
  &=\Tr(AX),\\
  \Tr\bigl[(D(z)^*AD(z)-A)X\bigr]&=0.
\end{align*}
This holds for every trace-class $X$.  Since trace class separates bounded
operators,
\begin{equation}\label{eq:A-invariant}
  D(z)^*AD(z)=A
  \qquad(z\in\C^s).
\end{equation}

\subsection*{Step 3: irreducibility forces a scalar}

Equation \eqref{eq:A-invariant} says that $A$ belongs to the Weyl commutant.
By \eqref{eq:weyl-commutant}, $A=cI$ for some scalar $c$.  The order bounds
$0\leq A\leq I$ imply $0\leq c\leq1$.  Substitution into
\eqref{eq:effect} yields
\[
  \Tr\Gamma(X)=\Tr(cIX)=c\,\Tr X,
\]
which is \eqref{eq:constant-trace}.

This also covers the endpoint cases.  If $c=1$, $\Gamma$ is trace
preserving.  If $c=0$, then $\Gamma(X)$ is a positive trace-class operator of
trace zero for every $X\geq0$, hence $\Gamma=0$.

\section{Intuition and examples}

The effect $A$ is the quantum analogue of a state-dependent acceptance
function.  Covariance says this function cannot distinguish any two points
connected by a phase-space displacement.  The representation is
irreducible, so there are no hidden invariant sectors on which different
acceptance probabilities could be assigned.  Therefore $A$ must be
constant.

For example, if $\mathcal A_\gamma$ is the quantum-limited amplifier, then
\[
  \Gamma=c\,\mathcal A_\gamma,
  \qquad 0\leq c\leq1,
\]
satisfies the lemma.  It performs the amplifier with an input-independent
probability $c$.  By contrast, a filter that accepts according to the vacuum
effect,
\[
  \Tr\Gamma(\rho)=\langle0|\rho|0\rangle,
\]
is not displacement covariant: displacing the state changes its vacuum
overlap.

Irreducibility is essential.  If a symmetry representation decomposed into
two invariant sectors, an effect of the form
$c_1P_1+c_2P_2$ would commute with the symmetry while assigning different
success probabilities to the two sectors.  Weyl displacements exclude this
possibility.

\section{Dependencies and role in the paper}

The proof uses only three standard facts:
\begin{enumerate}
  \item trace-class/bounded-operator duality, which gives
  $A=\Gamma^*(I)$;
  \item invariance and cyclicity of the trace;
  \item irreducibility of the finite-mode Weyl representation.
\end{enumerate}
It does not use the Gaussian fidelity bound or the stochastic-order lemmas
that follow it.

Lemma D.1 is used immediately in Lemma D.3.  There one starts with a
trace-nonincreasing covariant map.  Lemma D.1 writes it as
\[
  \Gamma=c\,\Gamma_0,
\]
where either $c=0$ and the claim is trivial, or $\Gamma_0=c^{-1}\Gamma$ is a
trace-preserving displacement-covariant amplifier.  The Gu\c{t}\u{a}--Matsumoto
least-noise theorem applies to $\Gamma_0$.  Multiplication by $c$ then turns
their normalized stochastic-order statement into the conic inequality needed
for a subnormalized output.  Lemma D.1 is therefore the bridge between
probabilistic limit maps and the usual theory of deterministic Gaussian
amplifiers.

\section{Minimal prerequisite checklist}

To read Lemma D.1, it is enough to know that a completely positive
trace-nonincreasing map represents a possibly probabilistic quantum
operation, that its success effect is $\Gamma^*(I)$, and that the multimode
Weyl displacement representation is irreducible.  The proof is precisely
Schur's lemma applied to the success effect.

\end{document}
